% ===================================================================
%  9. TAULA — Mendiak eta bainuetxea. — Basoa eta ehiza.
%  Traduction 2026 d'apres texto/io, controlee sur texto/fr, texto/en
%  et texto/es. Consigne : voir texto/eu/10-taula-01.tex.
%
%  LE TABLEAU DE LA MONTAGNE, ET LE BASQUE Y EST CHEZ LUI COMME A LA
%  FERME : « mendia » la montagne, « harkaitza » le rocher, « haitzuloa »
%  la grotte, « ibarra » la vallee, « iturburua » la source,
%  « leizea » le gouffre, « artzaina » le berger. Les mots du relief
%  sont les plus vieux de la langue, et ils survivent jusque dans les
%  noms de lieu du pays.
% ===================================================================

%%K t09-apar-1 apar
\VUsaut{4.0mm}
\VUtitre{15.0pt}{60}{9. TAULA}
\VUsaut{3.0mm}

%%K t09-tit-1 sub
\VUcentre{12.0pt}{0}{\textit{Lehen agerraldia.}}
\VUsaut{3.0mm}

%%K t09-tit-2 sub
\VUpk{11.4pt}{40}{Mendiak. --- Bainuetxeko herrixka.}
\VUsaut{3.0mm}

%%K t09-c1-01-1 p
1. --- Astebete daramat Pratzen. Ezin dut adierazi zenbateko harridura sentitu
nuen Pirinioetako \VUgras{mendikate} \textsuperscript{(1)} zoragarriaren aurrean
jarrita, zeinaren \VUgras{gailurra} \textsuperscript{(2)} zeru urdinean ebakitzen
baita, zerra erraldoi bat bezala.

%%K t09-c1-02-1 p
2. --- Ez nuen irudikatzen Pirinioetako \VUgras{tontorrak}
\textsuperscript{(3)} halako \VUgras{garaierara} irits zitezkeenik, eta harrituta
geratu nintzen \VUgras{glaziar} bikainekin \textsuperscript{(5)} eta
\VUgras{tontor} elurtuekin \textsuperscript{(4)}, zeinetatik hain
\VUgras{elur-jausi} beldurgarriak erortzen baitira.

%%K t09-tit-3 sub
\VUsaut{3.0mm}
\VUpk{10.2pt}{40}{Mendiko ikuspegia.}
\VUsaut{3.0mm}

%%K t09-c1-03-1 p
3. --- Iritsi eta hurrengo egunean bertan Pratzen ondoko \VUgras{sumendi}
itzali zaharrera \textsuperscript{(6)} joan nintzen. \VUgras{Kraterraren} ertz
biluzi eta antzuan oraindik argi ikusten dira \VUgras{labak} utzitako arrastoak,
zeina duela mende batzuk isuri baitzen \VUgras{mendiaren hegalean}
\textsuperscript{(7)} behera. \VUgras{Aldapetako} batean laba horren zati batzuk
aurkitzea lortu nuen.

%%K t09-c1-04-1 p
4. --- Pratz mugan dago, eta \VUgras{gotorlekua} \textsuperscript{(8)}, zeinak
Frantziaren eta Espainiaren arteko \VUgras{lepoa} \textsuperscript{(27)}
zaintzen baitu, oso antzekoa da Rhin ertzeko gaztelu zahar haiei, zeinek beren harkaitzaren
gainetik \VUgras{harrapakina} zelatatzen baitute.

%%K t09-c1-05-1 p
5. --- \VUgras{Arrano} erraldoiak \textsuperscript{(9)}, zeinaren habia
\VUgras{gotorlekutik} \textsuperscript{(8)} gertu baitago, sarritan erakartzen du
nire arreta. \VUgras{Harrapakin-hegazti} honek, \VUgras{moko} indartsuarekin,
sarritan ekartzen die txitoei bildots edo antxume bat, \VUgras{atzaparretan}
estututa. \VUgras{Hegoak} hain handiak ditu, non luzaroan hegan egin baitezake
bere zama astunarekin.

%%K t09-tit-4 sub
\VUsaut{3.0mm}
\VUpk{10.2pt}{40}{Oinezkoa.}
\VUsaut{3.0mm}

%%K t09-c1-06-1 p
6. --- Hemengo ibilaldirik atseginena Coo-ko \VUgras{ur-jauzira}
\textsuperscript{(10)} joatea da. \VUgras{Ur erorkorra} zurrunbiloan amiltzen da
eta gero \VUgras{erreka} polit bihurtuta desagertzen \VUgras{izei-basoan}
\textsuperscript{(11)}, herrixkatik mendebaldera. Baso horretatik igo ginen
\VUgras{meteorologia-etxolaraino} \textsuperscript{(12)} eta \VUgras{dorrearen}
gailurretik eskualdeko \VUgras{panorama} osoa hartu genuen begiz.
\VUgras{Ikuspegia} bikaina da, eta badirudi \VUgras{naturak} eskualde honi eman
dizkiola eskura zituen altxor guztiak.

%%K t09-c1-07-1 p
7. --- Jaisteko, \VUgras{funikularra} \textsuperscript{(13)} hartu genuen eta
\VUgras{geltoki} txiki batean jaitsi ginen, \VUgras{gaztelu feudal} zaharraren
\VUgras{hondakinen} \textsuperscript{(14)} ondoan, non garai batean Pratzeko
jaunak bizi baitziren.

%%K t09-c1-08-1 p
8. --- Gaztelu gaixoa! Zer pentsatu behar du bere azpian \VUgras{bainuetxeko
herrixka} \textsuperscript{(15)} txukuna ikusita, \VUgras{ur-bainuetxe} modan
bihurtua?

%%K t09-c1-08-2 p
\VUgras{Klima} hain da atsegina hemen, eta \VUgras{ur minerala} hain
sendagarria, non \VUgras{sendatzeak}, \VUgras{sanatorioan}
\textsuperscript{(17)} egiten denak, benetako atsegina ematen baitu.

%%K t09-tit-5 sub
\VUsaut{3.0mm}
\VUpk{10.2pt}{40}{Bainuetxeko herrixka atsegina.}
\VUsaut{3.0mm}

%%K t09-c1-09-1 p
9. --- Egia esan, dena egin da egonaldia atsegina izan dadin. \VUgras{Zubibide}
handi bat \textsuperscript{(16)} eraiki dute trenbiderako; \VUgras{kurtsaleko}
\VUgras{parkea} \textsuperscript{(18)}, non \VUgras{kontzertuak} ematen
baitira, xarmangarri diseinatuta dago. \VUgras{Txalet} dotore batzuk
\textsuperscript{(19)} eraiki dituzte \VUgras{kasinoaren} \textsuperscript{(21)}
inguruan, eta oso ondo zaindutako \VUgras{errepideak}
\textsuperscript{(22)} izugarri errazten du joan-etorria.

%%K t09-c1-10-1 p
10. --- \VUgras{Lubeta} soil batek \textsuperscript{(23)} bereizten du
\VUgras{bidea} \textsuperscript{(20)} \VUgras{ibarretik} bertatik
\textsuperscript{(25)}, zeinaren hondoan \VUgras{aintzira} liraina
\textsuperscript{(26)} baitago, \VUgras{erreka} garden bat \textsuperscript{(24)}
elikatzen duena.

%%K t09-c1-11-1 p
11. --- Ibarraren beste aldean burdinazko \VUgras{meategia}
\textsuperscript{(28)} ustiatzen dute. Hainbat aldiz joan naiz ikustera nola
jaisten diren \VUgras{meatzariak} \VUgras{putzura}, eta \VUgras{galerian} ere
sartu naiz, non eguna igarotzen baitute \VUgras{minerala} ateratzen, zeinetan
\VUgras{metala} baitago.

%%K t09-c1-12-1 p
12. --- Meategiaren ondoan \VUgras{urtzaile} handi bat eraiki dute, handik
ateratako aberastasunak bertan bertan erabiltzeko. \VUgras{Labe garaiak}
\textsuperscript{(29)}, hemengo \VUgras{ikatz} ugariaz berotuak, \VUgras{burdinurtua}
azkar eta merke lortzen uzten du.

%%K t09-c1-13-1 p
13. --- Meategitik gertu \VUgras{harrobia} \textsuperscript{(30)} ustiatzen dute.
\VUgras{Harrobiari} zaharrak bere \VUgras{aitzurraz} ez du \VUgras{granitorik}
askatzen, ez \VUgras{porfidorik}, ezta \VUgras{hareharririk} ere, baizik eta
etxeak eraikitzeko \VUgras{eraikuntza-harri} soila.

%%K t09-c1-14-1 p
14. --- Eskualde honen apaingarririk ederrenetako bat \VUgras{izeia}
\textsuperscript{(31)} da. Nonahi --- \VUgras{sakanaren} \textsuperscript{(32)}
hondoan, \VUgras{uharraren} \textsuperscript{(33)} ertzean, \VUgras{leizearen}
\textsuperscript{(34)} edo amildegiaren ondoan --- haren orratz finek beren nota
berdea ekartzen dute.

%%K t09-tit-6 sub
\VUsaut{3.0mm}
\VUpk{10.2pt}{40}{Oinezko ibilaldiak.}
\VUsaut{3.0mm}

%%K t09-c1-15-1 p
15. --- Oporrak baliatzen ditut ibilaldi luzeak egiteko. Menditik kiribiltzen
diren \VUgras{bidezidorrek} \textsuperscript{(35)} \VUgras{kilometro} batzuk
egiten uzten dute, eta sarritan \VUgras{lekoa} batzuk ere, distantzia ohartu
gabe.

%%K t09-c1-15-2 p
\VUgras{Mugarriek} \textsuperscript{(36)} eta \VUgras{seinaleek}
\textsuperscript{(37)} markatzen dituzte bidezidorrak, non sarritan topatzen
baita \VUgras{mendizale} gaztea \textsuperscript{(38)} bere \VUgras{zorroarekin}
\textsuperscript{(40)} saihetsean, \VUgras{marmota} bat \textsuperscript{(39)}
daramala, edo \VUgras{ijito} \textsuperscript{(41)} bat, zeinaren \VUgras{larruzko
txanoa} \textsuperscript{(42)} bitxiki eztabaidatzen baita \VUgras{kapa}
zahar urratuarekin \textsuperscript{(43)}. \VUgras{Hartz} ikasi bat
\textsuperscript{(44)} darama \VUgras{katez}.

%%K t09-tit-7 sub
\VUpk{10.2pt}{40}{Igoerak.}
\VUsaut{3.0mm}

%%K t09-c1-16-1 p
16. --- Ia egunero \VUgras{gidariarekin} \textsuperscript{(48)} joaten naiz
\VUgras{igoera} bat egitera, eta oso harro nago \VUgras{motxila}
\textsuperscript{(47)} bizkarrean eta \VUgras{alpenstocka} eskuan, benetako
\VUgras{turista} \textsuperscript{(46)} bat bezala, \VUgras{harkaitzei}
\textsuperscript{(45)} eraso egiten diedanean.

%%K t09-c1-17-1 p
17. --- Mendiaren gailurretik lautadako jendea inurriak baino handiagoa ez dela
dirudi; \VUgras{ardi-artaldeak} \textsuperscript{(50)}, \VUgras{larrean}
\textsuperscript{(49)} bazkatzen denak, haurren jostailu bat dirudi. \VUgras{Pinua}
\textsuperscript{(51)} edo zurezko \VUgras{txaleta} \textsuperscript{(52)} bera
ozta-ozta antzematen den puntu bat da berdetasunean.

%%K t09-c1-18-1 p
18. --- Ibilaldi hauetan sarritan topatzen dugu \VUgras{bidexkan}
\textsuperscript{(53)} \VUgras{sarrio-ehiztaria} \textsuperscript{(54)}, hil
berri duen animalia bizkarrean daramala.

%%K t09-c1-19-1 p
19. --- Nire herbariora \VUgras{iratzearen} \textsuperscript{(55)} hosto oso
aipagarri bat sartu dut, eta \VUgras{txilarraren} \textsuperscript{(57)}
adartxo arrosa polit bat, \VUgras{haitzulo} txiki baten \textsuperscript{(56)}
sarreran moztu nuena azken ibilaldian.

%%K t09-tit-8 sub
\VUsaut{3.0mm}
\VUcentre{12.0pt}{0}{\textit{Bigarren agerraldia.}}
\VUsaut{3.0mm}

%%K t09-tit-9 sub
\VUpk{11.4pt}{40}{Basoa. --- Ehiza.}
\VUsaut{3.0mm}

%%K t09-c2-01-1 p
1. --- Atzo egun zoragarria igaro nuen basoan. Han bizi diren animalien eta
\VUgras{intsektuen} \textsuperscript{(5)} bizitza oparoak asko interesatzen
ninduen.

%%K t09-c2-01-2 p
\VUgras{Armiarma} \textsuperscript{(1)}, bere \VUgras{sarea}
\textsuperscript{(2)} bi adarren artean ehuntzen, hegan igarotzen den
\VUgras{eulia} \textsuperscript{(3)} harrapatzeko; \VUgras{barraskiloa}
\textsuperscript{(4)}, bere etxetxoarekin astiro arrastaka; orkatz-kakalardoa,
harrapakina zelatan; inurri langilea, \VUgras{inurritegiaren}
\textsuperscript{(6)} ondoan lanean --- izaki txiki horiek guztiak ezohiko indarrez
zebiltzan eta lan egiten zuten.

%%K t09-c2-02-1 p
2. --- \VUgras{Sugea} \textsuperscript{(7)}, zeina lehenik \VUgras{sugegorritzat}
hartu bainuen, eta \VUgras{suge} kaltegabea zela ikusi ---, zuhaitz baten
enborraren azpian kiribildu zen eta noizean behin ateratzen zuen bere buru fina,
beti ziztatzeko prest bezala. Nire aurrean \VUgras{bare} handi bat
\textsuperscript{(8)} zihoan astiro, hemen ugariak diren \VUgras{marrubi} txiki
goxoetako bat \textsuperscript{(11)} jan zuena.

%%K t09-c2-03-1 p
3. --- Lurra \VUgras{basoko lorez} estalita zegoen: \VUgras{lore-berak}
\textsuperscript{(9)}, \VUgras{miruberak} \textsuperscript{(10)} eta ezagutzen ez
nituen beste landare asko loratzen ziren \VUgras{belar-mototsen}
\textsuperscript{(12)} artean.

%%K t09-c2-04-1 p
4. --- Eskualde honetan \VUgras{boilurra} ia \VUgras{perretxikoa}
\textsuperscript{(14)} bezain arrunta da, eta basozain baten \VUgras{txerrikumea}
\textsuperscript{(13)} gogotik ari zen arakatzen, aurkituko ote zituen.

%%K t09-tit-10 sub
\VUsaut{3.0mm}
\VUpk{10.2pt}{40}{Ehiza-lapurreta.}
\VUsaut{3.0mm}

%%K t09-c2-05-1 p
5. --- Asko dibertitu ninduen eszena baten \VUgras{amaiera} harrapatu nuen. Bere
\VUgras{lursagu-txakurraren} \textsuperscript{(15)} usaimenari esker
\VUgras{basozainak} \textsuperscript{(16)} \VUgras{ehiza-lapurra}
\textsuperscript{(22)} harrapatu berri zuen, hark hil zuen \VUgras{ehiza}
\textsuperscript{(17)} eramatera zihoan unean. Harrapakina uztea nahiago izanik
atxilotua ez izatearren, ehiza-lapurra ihes joan zen, eta \VUgras{erbia}
\textsuperscript{(18)}, \VUgras{oreina} \textsuperscript{(19)},
\VUgras{orkatza} \textsuperscript{(20)} eta \VUgras{basurdea}
\textsuperscript{(21)} belarraren gainean geratu ziren etzanda, basozainaren poz
handirako.

%%K t09-c2-06-1 p
6. --- Harrigarria da baso honetako zuhaitzen aniztasuna. \VUgras{Pagoek}
\textsuperscript{(23)} eta \VUgras{urkiek} \textsuperscript{(26)},
\VUgras{pinuek}, zeinek \VUgras{erretxina} ematen baitute
\textsuperscript{(27)}, \VUgras{haritzek} \textsuperscript{(29)} eta beste askok
elkarren artean gurutzatzen dituzte beren adarrak.

%%K t09-c2-06-2 p
\VUgras{Huntzak} \textsuperscript{(24)}, zuhaitz garaien enborrei itsasten
zaionak, berde bizia ematen dio \VUgras{baso garaiari} \textsuperscript{(25)},
zeina atsegin nahasten baita \VUgras{goroldioaren} \textsuperscript{(31)} tonu
zurbil eta epelarekin, non \VUgras{okil berdeak} \textsuperscript{(30)} intsektuak
bilatzen baititu.

%%K t09-tit-11 sub
\VUsaut{3.0mm}
\VUpk{10.2pt}{40}{Basoko ikuspegia.}
\VUsaut{3.0mm}

%%K t09-c2-07-1 p
7. --- Haritz zaharrek liluratu ninduten. Beren \VUgras{sustrai} korapilatsu
handiek \textsuperscript{(32)}, erdi lurretik aterata, indar eta ahalmen itxura
bikaina ematen diete, eta ez dut ulertzen nola erabaki dezakeen \VUgras{jabeak}
\textsuperscript{(35)} haiek eraistea eta \VUgras{zerrari}
\textsuperscript{(34)} ematea, \VUgras{egurgilearenari}
\textsuperscript{(33)}. \VUgras{Enbor} gaixoek \textsuperscript{(36)},
\VUgras{azalik} \textsuperscript{(37)} gabe utziak edo \VUgras{aizkoraz}
\textsuperscript{(38)} zatituak --- \VUgras{egun-langilearenaz}
\textsuperscript{(39)} ---, tristetu egiten naute.

%%K t09-c2-08-1 p
8. --- Ondoan \VUgras{ikazkin} zahar batek \textsuperscript{(41)} bere
\VUgras{palaz} ikatz \VUgras{pila} bat \textsuperscript{(42)} egin zuen, eta
emakume zahar batek \VUgras{txotx} \VUgras{sorta} bat \textsuperscript{(40)}
zeraman, moztutako zuhaitzen adarrez egina.

%%K t09-tit-12 sub
\VUsaut{3.0mm}
\VUpk{10.2pt}{40}{Ehiza.}
\VUsaut{3.0mm}

%%K t09-c2-09-1 p
9. --- \VUgras{Basozainaren etxean} \textsuperscript{(46)} pixka bat atseden hartu
nahi nuen, baina \VUgras{urmael} txikira \textsuperscript{(43)} iritsita, non
\VUgras{beltxarga} eder bat \textsuperscript{(45)} igerian baitzebilen, ohartu
nintzen \VUgras{uholde} berriak \textsuperscript{(44)} bidezidorra benetako
\VUgras{zingira} bihurtu zuela. Itzulinguru bat egin behar izan nuen, eta
basoertzean dauden \VUgras{zedroaren} \textsuperscript{(49)}, \VUgras{makalen}
\textsuperscript{(48)} eta \VUgras{zumarraren} \textsuperscript{(47)} ondotik
igarotzean, ikusi nuen nola \VUgras{sasitzatik} \textsuperscript{(54)} atera zen
\VUgras{orein} bikain bat \textsuperscript{(50)}, \VUgras{txakur-talde}
nekaezinak \textsuperscript{(51)} jarraitua, zeina \VUgras{adarrak}
\textsuperscript{(53)} berotzen baitzuen, \VUgras{ehiztari zaldunenak}
\textsuperscript{(52)}. Animalia gaixoa erabat akituta zegoen. Erori zen, hilzorian, niregandik
urrats gutxira.

%%K t09-c2-10-1 p
10. --- Basozainaren semea, \VUgras{ehizaren} zaratak erakarrita, etxetik atera
zen, eta biok basora itzuli ginen. \VUgras{Mika} bat \textsuperscript{(56)}
antzeman zuen haritz zahar baten adarrean. Uste zuen haren habia
\VUgras{hostoetan} \textsuperscript{(58)} ezkutatuta zegoela, eta hartu nahi zuen.

%%K t09-c2-10-2 p
\VUgras{Urtxintxa} \textsuperscript{(61)} bezain bizkorra, \VUgras{adarrez}
\textsuperscript{(55)} adar igotzen zen, baina ez zuen ezer aurkitu eta
\VUgras{bele} zahar bat \textsuperscript{(60)} bakarrik izutu zuen,
\VUgras{adaburuan} \textsuperscript{(57)} eserita zegoena.
\VUsaut{4.0mm}
\VUfilet{22mm}
